\documentclass[11pt]{article}

\usepackage[margin=25mm]{geometry}

\usepackage{amsmath,amssymb,fontspec,unicode-math}

\IfFontExistsTF{DejaVu Serif}{\setmainfont{DejaVu Serif}}{\setmainfont{Latin Modern Roman}}

\IfFontExistsTF{Latin Modern Math}{\setmathfont{Latin Modern Math}}{\setmathfont{DejaVu Math TeX Gyre}}

\usepackage{xurl,hyperref}

\setlength{\parindent}{0pt}\setlength{\parskip}{6pt}

\begin{document}

\title{Beyond Pairwise Connections: Extracting High-Order Functional Brain Network Structures under Global Constraints}

\author{ELVIA preliminary AI analysis}\date{21 October 2026}\maketitle

Preliminary assessment: Major.\par

Independent scientific review remains pending. Author adoption has not been requested.

The written covariance and conditional-variance arguments are false in their stated domains; the expressivity proof has separate gaps. Both released binarizers force complete graphs, and evaluation limitations require distinct known-label/unseen-row/window treatments.

\section*{Review scope}

Rapid current-version reassessment of every substantive concern in the supplied LaTeX report: the implicated paper passages and relevant author-released source paths were inspected, existing proposed corrections were briefly checked, and impact was calibrated to the demonstrated scope. No author-provided or historical verification program was executed.

\section*{Problem 1: Equal covariance does not make every bivariate-distribution edge rule equal}

Source: Theorem 1 and Corollary 1, PDF p. 5; Appendix C, PDF p. 25.

Assessment: supported; impact: major

\subsection*{Source claim (paraphrase)}

Claim and location (paraphrase). Theorem 1 and Corollary 1, p.5, consider dimension at least three, equality of all second-order lagged moments used by the edge rule, and at least one differing cumulant of order at least three. They conclude that every pairwise generator from the preceding bivariate-distribution definition returns the same adjacency. Appendix C, p.25, supplies the proof.

\subsection*{Analysis}

At every time independently, process \(A\) is \((U,V,UV)\) for independent fair Rademacher variables \(U,V\). Process \(B\) has independent coordinates each taking \(-\sqrt2,0,\sqrt2\) with probabilities \((1/4,1/2,1/4)\). Both are bounded, stationary, centered, have covariance \(I_3\), and have zero nonzero-lag covariances. Their third joint cumulants are \(1\) and \(0\).

The permitted bivariate edge functional \(\mathbb E[X_1^4X_2^2]\) is \(1\) for \(A\) and \(2\) for \(B\). Thresholding it at \(3/2\) gives different binary edges. Equal second moments therefore do not make every bivariate-distribution edge rule equal. Restricting the rule to covariance, or matching full bivariate laws, restores the corresponding restricted implication.

\subsection*{Finding and scope}

The theorem's universal bivariate-rule conclusion is false. This narrows its impossibility claim while preserving the motivation that some higher-order interactions are invisible to pairwise information.

\subsection*{Proposed correction}

Restrict the rule to matched second-order inputs, or strengthen the premise to equality of the full bivariate laws used by each rule. Equal-pair-law XOR examples still establish a restricted higher-order non-identifiability result.

A local correction or qualification is proposed where specified; author adoption is not established.

\section*{Problem 2: Positive conditional variance does not imply different marginal distributions}

Source: Appendix B, Lemma 2, PDF pp. 24–25; Theorem 3.

Assessment: supported; impact: major

\subsection*{Source claim (paraphrase)}

Claim and location (paraphrase). Appendix B, pp. 24–25, Lemma 2 states that \(\operatorname{Var}(X\mid Y)>0\) almost surely implies \(X\) and \(Y\) are not equal in distribution. Theorem 3 applies this along project, group, subject, and sample levels and concludes all six pairs of graph distributions differ.

\subsection*{Analysis}

Independent \(X,Y\sim\operatorname{Bernoulli}(1/2)\) have the same marginal law, yet \(\operatorname{Var}(X\mid Y)=1/4>0\) almost surely. Equal distributions do not imply equality on each outcome. Likewise, take four independent Bernoulli edge indicators for two-node graphs at the hierarchy levels. Every parent-conditioned edge variance is \(1/4\), but all four graph laws coincide. Neither the lemma nor all six claimed pairwise distinctions follow.

\subsection*{Finding and scope}

The formal distributional-distinction argument is false. Different semantic hierarchy levels or empirically different fitted graphs do not repair that universal implication.

\subsection*{Proposed correction}

State conditional nondeterminism as the consequence of positive conditional variance; provide assumptions or statistics distinguishing marginal laws before asserting distributional separation.

A local correction or qualification is proposed where specified; author adoption is not established.

\section*{Problem 3: The higher-order expressivity proof contains independent invalid steps and does not prove optimization convergence}

Source: Theorem 2, PDF p. 6; Section 7, PDF p. 10; Appendix D, Theorem 4 and Lemmas 3–5, PDF pp. 26–27.

Assessment: supported; impact: major

\subsection*{Source claim (paraphrase)}

Claim/location (paraphrase). Theorem 2, p.6, and Theorem 4, Appendix D, pp. 26–27, claim arbitrarily small classification risk for measurable higher-order rules and a class-separable measurable label map. The sentence following Theorem 2, and Section 7 on p.10, additionally attribute convergence of GCM to that theorem. Appendix D uses Lemmas 3–5.

\subsection*{Analysis}

The printed sampler first places its edge parameter \(A\) in \((0,1)\). For standard Gumbel noise \(M\), the desired off-edge event satisfies \(\Pr(\operatorname{sigmoid}((A+M)/\tau)<1/4)=\exp[-\exp(A+\tau\log3)]\le e^{-1}<0.368\). Increasing the earlier raw parameter cannot make that event arbitrarily probable after the intervening sigmoid.

A separate lemma claims any continuous injective node map has an injective tagged sum. The identity map with distinct one-hot tags gives equal sums for feature triples \((0,2,0)\) and \((1,1,0)\), despite different tagged inputs. Finally, a continuous MLP cannot uniformly approximate the measurable step \(\mathbf1\{x_1\ge0\}\) below error \(1/2\) on the whole cube. These invalidate the supplied expressivity argument; representational existence would not establish optimizer convergence in any event.

\subsection*{Finding and scope}

The supplied proof establishes neither its unrestricted expressivity statement nor optimization convergence. Suitable restricted or alternative constructions may still be possible; this is not a claim that higher-order learning is impossible.

\subsection*{Proposed correction}

A revised proof must use a sampler capable of the stated graph approximation, an explicitly suitable tagged encoder, and a valid probability/risk approximation argument. Treat optimizer convergence separately from representational existence.

A local correction or qualification is proposed where specified; author adoption is not established.

\section*{Problem 4: Both released binarizers return all-ones graphs; the released loss also omits the stated sparsity term}

Source: Section 5, Eq. (3), PDF p. 6; BBA ablation, PDF pp. 8–9; Algorithm 2, PDF p. 27; current models.py and official supplementary code/models.py.

Assessment: supported; impact: minor

\subsection*{Source claim (paraphrase)}

Claim/location (paraphrase). Section 5, p.6, Eq. (3), Algorithm 2, p.27, and the BBA ablation discussion, pp.8–\(9\), require selection of a bounded number of high-scoring edges. The method’s loss is \(L=L_1+\alpha L_2+\beta R\) with an \(\ell_1\) sparsity term. The relevant released functions are in the current models.py and in the official supplemental code/models.py.

\subsection*{Analysis}

Both released binarizers scatter scalar \(1\) into sorted\_indices, a permutation of all \(N^2\) matrix positions, while leaving the computed top-k mask unused. Every forward adjacency is therefore the all-ones matrix, including its diagonal. A straight-through estimator can still return a nonzero surrogate gradient, which does not change that constant forward graph. The active loss separately omits the stated sparsity term. The historical table/figure-producing program remains unidentified.

\subsection*{Finding and scope}

The release cannot produce sample-dependent sparse hard graphs through these functions. This is a substantial release defect, but published-figure provenance and central empirical damage are not established, so its impact label is Minor.

\subsection*{Proposed correction}

For the intended undirected loop-free forward graph, select a rank mask on strict-upper-triangle scores, mirror it and zero the diagonal. Explicitly define the edge budget, backward estimator and actual regularization objective; a one-line scatter edit alone is incomplete.

A local correction or qualification is proposed where specified; author adoption is not established.

\section*{Problem 5: Evaluation requires explicit separation of known-label group analysis, unseen graph parameters, and possible duplicate windows}

Source: Section 6.2, PDF p. 8; Section 8, PDF p. 10; released graph lookup and window-split construction.

Assessment: supported; impact: minor

\subsection*{Source claim (paraphrase)}

The source distinguishes group-conditioned structure analysis from held-out classification, while the release allocates graph rows and selected window identities before its split.

\subsection*{Analysis}

Three evaluation scopes differ. True-label group routing is explicitly acknowledged by the authors and supports class-conditioned exploration. Held-out sample/subject lookup rows receive no training gradient without an inductive graph generator. With-replacement window selection can duplicate one original window under different new IDs on opposite split sides. The all-ones bug masks some graph differences, and no published leakage percentage is measured.

\subsection*{Finding and scope}

Known-label routing is a preserved author-acknowledged clarification; held-out-row interpretation and possible duplicate-window leakage are bounded additional release concerns.

\subsection*{Proposed correction}

Keep label-conditioned group structure exploration separate from label-blind classification. Specify an inductive or adaptation protocol for unseen graph rows and split by preserved original-window identities to prevent duplicate-window overlap.

A local correction or qualification is proposed where specified; author adoption is not established.

\section*{Limits}

No author-provided program, training procedure, model inference or wet-lab experiment was executed in this review.\par

Historical execution results are retained only with that attribution; a source-level defect alone does not identify the program or configuration used for the published experiments.\par

No author contact, adoption, independent human scientific review or novelty determination is established.\par

The findings apply only to the identified claims, source versions and stated assumptions; they do not certify or invalidate the entire paper.\par

\section*{Source identity}

arXiv:2510.09175v1 (10 October 2025)

\url{https://arxiv.org/pdf/2510.09175v1}

Reviewed PDF SHA-256: \nolinkurl{343d9565d687c3e5ef2870d31bef55c4116818b42d73c0302e2253cb1d1b9428}

The exact inspected current paper PDF is identified by its SHA-256. The current arXiv paper matches all five authors and retains the formal paper’s implicated definitions and proofs at the cited locations. Current author head c42ed5cdb80962db871a62ad4474921cc10fc0c8 retains the all-position scatter and evaluation routes; official supplementary code was separately read. An exact PDF acquisition timestamp is not established in this reassessment; source-currentness checks have their own recorded timestamps and file modification times were not substituted.

\end{document}
